\section{Uniform half density}\label{sec:uniform-density}

We first give a construction that retains an arbitrary $C_4$-free base.
The only interaction condition concerns a matching of columns: no fixed
row may contain two of its edges. A rainbow matching will supply this
condition in the asymptotic argument.

\begin{lemma}[Column-fan gluing]\label{lem:column-fans}
Let $J$ be a matching of columns of a $C_4$-free base $F$, such that each
row of $F$ contains at most one edge of $J$. For $ab\in J$, let $H_{ab}$
be the rows missing both $a,b$, and suppose $|H_{ab}|\ge2$.
There is at most one row $z$ fixed in both columns. If it exists, color
each row of $H_{ab}$ by its common fixed column with $z$, giving separate
colors to rows with no such column. Suppose $H_{ab}$ can be ordered
without equal consecutive colors. If no $z$ exists, use any ordering.
Then
\[
 R(F)\ge\sum_{ab\in J}(|H_{ab}|-1).
\]
Here $R(F)$ denotes the pair optimum on this base even when $F$ is not maximum.
\end{lemma}
\begin{proof}
Write an ordering as $q_0,\ldots,q_L$ and select
$A_t=\{(q_t,a),(q_{t+1},b)\}$ for $0\le t<L$.
The two endpoint cells $(q_0,b),(q_L,a)$ are holes. Different fans use
disjoint column sets. First normalize every chain backwards: the terminal
$a$-hole starts the induction. Once $A_{t+1}$ is identified, its representative
$(q_{t+2},b)$ shares column $b$ with the other companion $(q_t,b)$
of $A_t$. This is a column ground, or a hole ground if that companion
is absent. No product proved later is used.

For a foreign column, let $Y_k$ denote its coefficient at $q_k$.
A row grounds $A_tY_k$ when $k=t,t+1$. When $k<t$, an $a$-rectangle
transfers the target to $A_kY_t$. When $k>t+1$, a $b$-rectangle gives
$A_{k-1}Y_{t+1}$; this is a row ground for $k=t+2$, and otherwise an
$a$-rectangle gives $A_{t+1}Y_{k-1}$. The gap decreases by two.
Thus every pair is orthogonal to every foreign coefficient on its path.

Pairs of one fan share a column. For a pair $P$ from another fan, the
previous paragraph applies if one of its rows lies on this path. Otherwise
both rows of $P$ lie outside it, where columns $a,b$ contain only fixed
cells or holes. $C_4$-freeness supplies a hole in that two-by-two block;
choosing the corresponding representatives grounds the product. All
pair/pair products are therefore proved.

Consider a fixed cell $(y,d)$. Path rows and own fan columns are already
covered. If an exterior row $y$ misses $a$ or $b$, that cell is still a
genuine hole: other fans use other columns. It directly grounds the target.
Only the common fixed row $z$ remains. For a fixed $(z,d)$ outside $a,b$,
the ordering gives an endpoint row $r$ of the pair with $(r,d)\notin F$.
Use the rectangle through this endpoint and $(z,d)$. Its companion at
$(r,d)$ is a hole or a pair $Q$ in another fan. In the latter case $z$
misses that fan's other column, since $z$ contains at most one edge of $J$.
Since $z$ is fixed in $d$, it lies outside that fan's path. Its missing
cell in the other column is therefore a genuine hole. Use $Q$'s
representative in that column---available because all fans were normalized
first---to ground its product with $(z,a)$ or $(z,b)$. This is a direct
hole ground, so the common-root cases introduce no circular dependency.

Finally, a rectangle between two fixed cells has a nonfixed companion
by $C_4$-freeness. Its deviation is zero by a hole or a pair/fixed or pair/pair product
already proved. The two companions cannot be one selected pair: that
pair's opposite corners lie on rows missing both of its columns and
therefore cannot be the original two fixed cells. Line products are
immediate. These are all the recursive-line rules, and every singleton
has been retained.
\end{proof}

The coloring in the lemma is well-defined by $C_4$-freeness. A multiset
can be ordered without equal consecutive colors whenever its largest
color class has size at most half the total, rounded up: list the colors
in decreasing multiplicity and fill alternate positions first, then the
remaining positions. Thus the lemma
applies if every column incident with the common row $z$ has degree
$o(|H_{ab}|)$, uniformly over the selected edges.

\subsection{Choosing column pairs}

We use the following matching lemma.
Its proof uses Tutte's theorem and the de Bruijn--Erd\H{o}s incidence theorem.
\Needspace{7\baselineskip}
\begin{lemma}[Packing blocks in a linear set system]\label{lem:block-capacity}
Let a set of $N$ points be partitioned into parts $\mathcal P$, and let
$\mathcal L$ be a family of subsets of size at least two. Distinct members
of $\mathcal L$ meet in at most one point, and each meets every part in
at most one point. For $q\ge1$, there are $q$ disjoint pairs, none inside a part and
at most one inside each member of $\mathcal L$, if and only if
\[
 N\ge2q,\qquad |P|\le N-q\ (P\in\mathcal P),\qquad
 |L|\le N-q+1\ (L\in\mathcal L).
\]
\end{lemma}
\begin{proof}
Necessity follows by counting points outside a part or a member of
$\mathcal L$. First suppose $N=2q$. For each $L$, add two adjacent
vertices and join both to every point of $L$. Join two original points
directly precisely when they lie in different parts and in no common
member of $\mathcal L$. A perfect matching of this graph gives the
required pairs: the two vertices of $L$ either match each other or
match two points of $L$. The intersection assumptions make these
assignments unambiguous.

Apply Tutte's theorem \cite{Tutte}. Delete a vertex set $S$, containing
$h$ original points, both added vertices of $r$ members of $\mathcal L$,
and one added vertex of $s$ others. Thus $|S|=h+2r+s$. Components with
no original point contribute at most $s$ odd components. If the number
$c$ of components containing points is at most $h+2r$, Tutte's inequality
follows. Assume $c>h+2r$.

If $c=1$, then $h=r=0$. For $s>0$, every remaining added vertex still
has a point neighbor, so there is at most one odd component, at most
$|S|$. For $s=0$, the graph has even order and its only point-containing
component is even. If $r=0$ and $c\ge2$, any two points in different
components must belong to the same part: a common member of
$\mathcal L$ would still have an added vertex joining them. All remaining points therefore
belong to one part, giving $c\le2q-h\le q\le h$, a contradiction.

It remains that $r\ge1$ and $c\ge3$. Choose one point from each component.
Every pair of representatives belongs to exactly one part or one of
the $r$ removed members of $\mathcal L$. If the representatives are not
all in one such set, the de Bruijn--Erd\H{o}s incidence theorem
\cite{deBruijnErdos} requires at least $c$ sets. There are at most
$r+\lfloor c/2\rfloor$, since the parts are disjoint. This contradicts
$c>h+2r$. Thus every choice of representatives is contained in one set.
Keeping two representatives and varying a third shows that all remaining
points lie in that same set. A part gives $c\le h$ as above. A member of
$\mathcal L$ gives $2q-h\le q+1$, hence
$c\le2q-h\le h+2\le h+2r$, again a contradiction. Tutte's theorem applies.

For general $N\ge2q$, add $N-2q$ dummy points as one new part, meeting
no member of $\mathcal L$. There are $2(N-q)$ points, and the balanced
case applies with $q$ replaced by $N-q$. Delete the pairs meeting
dummy points. The remaining $q$ pairs have the required properties.
\end{proof}



\subsection{Below a fixed quadratic scale}

\begin{proof}[The first range in Theorem~\ref{thm:uniform-density}]
Orient $m\ge n$, put $M=|F|=z(m,n)$, and fix
$\varepsilon=1/100$, $c_0=\varepsilon^2$.
For all sufficiently large $m$, uniformly over $1\le n\le m$,
\begin{equation}\label{eq:ordinary-lower-uniform}
 M\ge(1-\varepsilon)n\sqrt m.
\end{equation}
Indeed choose a prime $t=(1-o(1))\sqrt m$ with $t^2+t+1\le m$,
and keep $\min(n,t^2+t+1)$ columns of its incidence matrix.
The ratio of this support's size to $n\sqrt m$ tends to one uniformly:
if the plane has fewer than $n$ columns, use $n\le m$.
The prime number theorem and incidence construction suffice
\cite{Singer,Hadamard}.

Suppose $n\le m\le c_0n^2$ and $n$ is sufficiently large.
Pair counting gives $M\le n\sqrt m+m\le2\varepsilon n^2$.
For row degrees $d_r$ and $\mu=M/m$, using
\eqref{eq:ordinary-lower-uniform},
\[
 V:=\sum_r(d_r-\mu)^2
 \le n(n-1)+M-\frac{M^2}{m}\le4\varepsilon n^2,
 \qquad \mu\le\sqrt n+1\le2\varepsilon n.
\]
Consequently $\max_r d_r\le(2\varepsilon+2\sqrt\varepsilon)n<n/2$.
Apply Lemma~\ref{lem:block-capacity} to the $n$ columns, with singleton
parts and with fixed-row neighborhoods of size at least two as
$\mathcal L$. These sets are distinct and meet pairwise in at most one
column, by $C_4$-freeness. With $q=\lfloor n/2\rfloor$, all capacity
conditions hold, including $d_r\le n-q+1$.
We obtain $q$ disjoint column pairs with at most one inside each fixed row.

A column of degree $E$ has at most $E+n-1$ fixed cells in its incident
rows. Applying pair counting to the remaining rows gives
$M\le m+n-1+n\sqrt{m-E}$. Therefore
\[
 \sqrt{1-E/m}\ge1-\varepsilon-\frac{\sqrt m}{n}
                         -\frac1{\sqrt m}\ge1-3\varepsilon,
 \qquad E\le6\varepsilon m<m/5.
\]
For every matched pair $ab$, its common-hole set has at least $3m/5$
rows, and every ordering color class has at most $m/5$ rows (including
fresh singleton classes for large $m$). Lemma~\ref{lem:column-fans} applies.
Writing $e_a$ for a column degree, disjointness of the matched columns gives
\begin{equation}\label{eq:fan-exact-sum}
 M+R(F)\ge M+\sum_{ab\in J}(|H_{ab}|-1)
          \ge |J|(m-1).
\end{equation}
Here $|H_{ab}|=m-e_a-e_b+|N_F(a)\cap N_F(b)|$, and the sum of all
matched column degrees is at most $M$. Thus
$M+R(F)\ge\lfloor n/2\rfloor(m-1)\ge mn/2-m$.
\end{proof}

\subsection{The quadratic and larger ranges}

The following elementary replacement bounds row degrees when the
number of rows is at least a fixed multiple of $n^2$.

\begin{lemma}\label{lem:bounded-row-trade}
For each $c>0$, the row degrees of every maximum $m\times n$ base
with $m\ge cn^2$ are bounded by a constant depending only on $c$,
for all sufficiently large $n$.
\end{lemma}
\begin{proof}
Pair counting gives $M/m\le1+1/\sqrt c$. Choose an integer
$K\ge2(1+1/\sqrt c)$. At least $m/2\ge n$ rows have degree at most
$K$, for sufficiently large $n$; these differ from the high row used below.
Suppose a row has degree $a>2\binom{K+1}{2}^2+1$ and put $k=K+1$.
Remove that row and $a$ low-degree rows. Let $A$ be its $a$ columns;
every surviving row meets $A$ at most once.

There are $a+1$ $k$-subsets of $A$ with no repeated pair. Construct
them greedily: after at most $a$ choices there are at most
$a\binom k2$ forbidden pairs, and a uniformly chosen $k$-subset
contains fewer than one of them in expectation, since
\[
 \frac{a\binom k2^2}{\binom a2}<1.
\]
It therefore avoids every forbidden pair for some choice.
Insert these $a+1$ blocks in the freed rows. The result is $C_4$-free.
At most $a+aK=ak$ cells were removed, but $(a+1)k$ were inserted,
contradicting the maximum cardinality of $F$.
\end{proof}

\begin{proof}[Completion of Theorem~\ref{thm:uniform-density}]
First assume $m\ge cn^2$ for fixed $c>0$ and sufficiently large $n$.
By Lemma~\ref{lem:bounded-row-trade}, choose an integer $D\ge2$,
with $D=O_c(1)$, bounding all row degrees. Color each column pair by its unique common fixed row,
or give it a fresh color when none exists. Set $\varepsilon=1/5$. Call a column heavy if
its degree exceeds $\varepsilon m$. Since $M=O_c(m)$, there are
$O_c(1)$ heavy columns. Omit these vertices from the column graph.
Mark as forbidden each edge whose common fixed row meets a heavy column.
For one heavy column $h$, $C_4$-freeness gives
\[
 \sum_{z\sim h}\binom{d_z-1}{2}
 \le\frac{D-2}{2}\sum_{z\sim h}(d_z-1)
 \le\frac{(D-2)(n-1)}2.
\]
Thus only $O_c(n)$ edges are forbidden.

Omit one further vertex if necessary so the number $N=n-O_c(1)$
of remaining columns is even. This parity deletion does not enlarge the set
of forbidden row colors. Choose a uniform perfect matching of the complete
graph on these vertices, with $N\ge4$. A given edge is selected with
probability $1/(N-1)$, so the expected number of forbidden matching edges
is $O_c(1)$. Two disjoint edges are selected together with probability
$1/((N-1)(N-3))$. Each row color has at most $B=\binom D2$ edges, and
\[
 \sum_{\text{row colors }i}\binom{b_i}{2}
 \le\frac{B-1}{2}\binom n2=O_c(n^2).
\]
Thus the expected number of equal-color pairs of matching edges is $O_c(1)$;
intersecting edges cannot both be selected, and fresh colors never collide.
Some matching has total cost $O_c(1)$. Delete its forbidden edges
and all but one edge of each repeated color. A color appearing $r$ times
costs at most $r-1\le\binom r2$ deletions. We obtain a rainbow matching
of $n/2-O_c(1)$ edges.

For each retained matching edge $ab$, both columns have degree at most
$m/5$, and their common fixed row, if any, meets no heavy column.
Hence $|H_{ab}|\ge3m/5$, while every ordering color class has size at
most $m/5$. Lemma~\ref{lem:column-fans} and the exact summation
\eqref{eq:fan-exact-sum} give
\begin{equation}\label{eq:quadratic-error}
 M+R(F)\ge |J|(m-1)\ge\frac{mn}{2}-O_c(m).
\end{equation}
The constants depend only on $c$, uniformly in $m\ge cn^2$ and in $F$.

Take the fixed $c=c_0=10^{-4}$ from the first range. Together, the two
arguments give $M+R(F)\ge mn/2-Cm$ for some absolute $C$, for all
$m\ge n\ge n_0$. Increasing $C$ to at least $n_0/2$ covers $n<n_0$
as well, since the claimed lower bound is then nonpositive.
Every line, rectangle and normalization rule is invariant under exchanging
rows and columns, so transposition gives the stated bound for all dimensions.

The cell caps $S(F)\le(mn-M)/2$ and $\zt(m,n)\le(mn+M)/2$, together
with $\zr(m,n)\ge M+R(F)$, give both comparison bounds.
Finally $M/(mn)\le1/n+1/\sqrt m\to0$ when both dimensions grow.
The lower bound just proved and the cell cap therefore give the uniform
half-density limits.
\end{proof}

